\section{Introduction}

\lettrine[lines=2]{\color{titleblue}W}{ ater} distribution networks are critical infrastructures for urban environments, ensuring the continuous delivery of potable water to residential, industrial, and commercial users. However, these systems are highly susceptible to leaks due to aging infrastructure, pressure fluctuations, and environmental factors. Undetected leaks can lead to significant water losses, increased operational costs, reduced service efficiency, and long-term structural damage. Consequently, the development of reliable and automated leak detection systems has become an important research and industrial challenge.

Traditional leak detection methods typically rely on manual inspection, acoustic sensing, or rule-based thresholding techniques applied to pressure and flow measurements. While these approaches can be effective in controlled scenarios, they often lack scalability and adaptability in complex urban networks. Moreover, such methods struggle to capture the temporal dynamics and nonlinear patterns associated with leak events, particularly in the presence of noise and varying operating conditions.

With the increasing deployment of sensor networks and Internet of Things (IoT) technologies in smart water systems, large volumes of multivariate time-series data are now available. This has opened the door for data-driven approaches, particularly deep learning methods, which have demonstrated strong capabilities in modeling sequential data. Architectures such as Convolutional Neural Networks (CNN) and Long Short-Term Memory (LSTM) networks have been widely used for time-series classification tasks due to their ability to capture local and long-range dependencies, respectively. However, existing approaches often rely on single-model architectures and lack comprehensive evaluation against alternative methods, limiting their ability to fully exploit complementary modeling strengths.

In this work, we propose a hybrid deep learning architecture that combines CNN, Bidirectional LSTM (BiLSTM), and a self-attention mechanism for automatic water leak detection. The CNN component extracts local temporal features from sensor signals, while the BiLSTM captures bidirectional dependencies across time. The integration of a self-attention layer allows the model to focus on the most informative segments of the input sequence, improving both detection performance and interpretability.

Beyond the model architecture, we design a complete end-to-end pipeline tailored for real-world deployment. This includes domain-informed feature engineering, sliding-window sequence construction, and a strict temporal separation between training, calibration, and test sets to prevent data leakage. Furthermore, we address the inherent class imbalance in leak detection by adopting evaluation metrics such as AUC-PR and by performing threshold calibration to control the trade-off between detection sensitivity and false alarm rate.

To ensure a rigorous evaluation, the proposed model is compared against several baseline machine learning and deep learning approaches. This benchmark analysis provides a comprehensive assessment of the effectiveness of the hybrid architecture relative to alternative methods.

Experimental results on a real-world dataset demonstrate that the proposed approach achieves strong performance, with high AUC-ROC and AUC-PR scores, as well as a favorable balance between precision and recall. In particular, the model maintains a low false alarm rate while achieving high leak detection accuracy, making it suitable for practical monitoring systems where both reliability and responsiveness are essential.

The main contributions of this paper can be summarized as follows:
\begin{itemize}
    \item A hybrid CNN-BiLSTM architecture enhanced with a self-attention mechanism for multivariate time-series leak detection.
    \item A robust data processing pipeline including feature engineering, sliding-window transformation, and leakage-free temporal splitting.
    \item An imbalance-aware evaluation framework incorporating AUC-PR and threshold calibration for operational decision-making.
    \item A comprehensive benchmark comparison against classical and deep learning models to validate the effectiveness of the proposed approach.
\end{itemize}